\chapter{Discussion and Future Directions}
\label{ch:discussion}


%\chapter{Discussion}

%\section{Integrating Theory into Practice}

%\subsection{Evaluating SpeechPrint: Demand, Limitations, and Outcomes}

%\subsection{Independent Practice and Learner Implications}

%\subsection{SpeechPrint’s Limitations and Future Directions}

%\section{Conclusions}

%\chapter{Future Horizons: From Annotation to Speech Authoring}
%\label{ch:future}

%The work presented in this thesis demonstrates that the fundamental bottleneck in modern speech technology is not the \textbf{extraction} of signal, but its \textbf{representation}. While current Automatic Speech Recognition (ASR) systems perform a form of "lossy compression" by flattening the multidimensional human voice into a one-dimensional lexical stream \cite{krifkadobes2026}, \textbf{SpeechPrint} serves as a Universal Metadata Layer that recovers the discarded colors and textures of human intent. Looking forward, the evolution of this symbolic layer points toward a transition from passive analysis to active authoring, grounded in \textbf{vowel-centric reasoning} and \textbf{verifiable alignment}.

%\section{The Vowel-Nuclei Paradigm and Symbolic Grounding}
%A critical advancement for the next iteration of SpeechPrint is the full adoption of the \textit{Carrier Hypothesis}. Recent research in \textit{VowelPrompt} (2026) asserts that vowels serve as the primary stable carriers of affective prosody, hosting the most salient pitch and energy patterns while consonantal regions remain spectrally unstable \cite{wang2026vowelprompt}. 

%By narrowing "etic" extraction specifically to vowel nuclei, future versions of the SpeechPrint pipeline can achieve a significantly higher signal-to-noise ratio for emotion recognition and prosodic highlighting. This allows the symbolic layer (/, \textbackslash, –, *) to move beyond mere annotation and into the realm of \textbf{paralinguistic-aware Spoken Language Models (SLMs)}. Augmenting transcripts with fine-grained, vowel-level natural language descriptors enables Large Language Models to "hear" emotions from text by reasoning over pitch slope and intensity variation rather than consuming opaque, uninterpretable acoustic embeddings \cite{wang2026vowelprompt, wu2025beyond}.

%\section{The Speech DAW: Disentanglement and Active Authoring}
%The current prototype of SpeechPrint facilitates a "literacoustic" reading of text, but its ultimate potential lies in transitioning into a \textbf{Speech DAW (Digital Audio Workstation)} for authoring intent. Traditional tools separate text (ELAN) and prosody (Praat), whereas SpeechPrint investigates how these can be reintroduced through a unified representation \cite{krifkadobes2026, paschen2020doreco}.

%By adopting disentangled speech architectures—such as those pioneered by the \textit{Dis-Vector} project—speech can be decomposed into independent "latents" for content, pitch, rhythm, and timbre \cite{disvector2026}. In this future workstation, a user—whether a poet, voice actor, or educator—could click a specific word and "re-pitch" or "re-rhythm" the output using SpeechPrint's symbolic layer as the control surface. This treats speech as \textbf{Dynamic Vector Graphics} not just for scientific visualization, but for zero-shot speech editing and style transfer where phonetic units become navigable, editable objects \cite{pataca2021speech, mccurdy2016poemage}.

%\section{Verifiable Intent and Clinical Impact}
%As generative AI approaches "perfect" prosodic output, the role of a symbolic representation becomes one of \textbf{verification and reasoning} rather than just generation. Future developments should leverage Reinforcement Learning with Verifiable Reward (RLVR) to ensure that the "sonic topology" generated by an AI aligns precisely with a human researcher's goals \cite{deepseek2025r1, wang2026vowelprompt}.

%This has profound implications for \textbf{Oral Reading Fluency (ORF)} and clinical linguistics. It allows for the automated scoring of student expressiveness against established benchmarks, such as the NAEP Fluency Scale, providing immediate and interpretable feedback that standard "flat" ASR cannot offer \cite{sammit2022automated, naep2005}. By bridging the gap between raw audio and human-readable metadata, SpeechPrint provides the transparent interface required for humans to remain the primary authors of their own vocal intent.

%\section{Conclusion: Preservation as Originality}
%The "SpeechPrint" philosophy argues that \textbf{Speech $\neq$ Text}. The originality of this work lies in the refusal to accept the information loss inherent in standard transcription. By elevating prosody to a first-class annotation layer, we preserve the "phonetic fingerprint" of identity and the sonic topology of human art \cite{krifkadobes2026, clement2013sounding}. In an era of increasingly black-box AI, the symbolic grounding of prosody ensures that the "voice" of the speaker survives the transition to the digital page.



\newpage



% ADDED HERE — future directions: TTS resynthesis, prosody DAW, classification model
% (from discussion with Lonce: he mentioned TTS and prosody control; also my own idea for a Speech DAW)

\section{Future Directions}
\label{sec:future}

\subsection{Prosody Resynthesis: From Annotation to Sound}
\label{sec:resynthesis}

A natural extension of the SpeechPrint pipeline is to close the loop between
analysis and synthesis. The current system moves in one direction: from audio
to symbolic annotation. The inverse direction -- from symbolic annotation to
audio -- would constitute a \emph{prosody resynthesis} capability. This concept
is illustrated in Figure~\ref{fig:resynthesis_pipeline}.

The proposed pipeline would work as follows. First, a recording is processed by
the existing SpeechPrint pipeline to produce a TextGrid with a prosody tier.
Second, a user modifies the prosody symbols in the TextGrid, for example,
changing a falling accent (\texttt{*\textbackslash}) to a rising one
(\texttt{*/}) on the word \textit{Mary} to produce the reading ``MARY flew to
Milan yesterday'' rather than ``Mary FLEW to Milan yesterday.'' Third, the
modified symbolic annotation is translated into prosody control parameters that
are passed to a speech synthesis system. The synthesis system produces a new
audio file with the requested prosodic pattern while preserving the voice quality
and segmental content of the original.

The linguistic term for this operation is \emph{prosody transfer} or
\emph{intonation resynthesis}. It has been studied in the context of
human-to-human style transfer and as a component of expressive text-to-speech
synthesis \cite{skerry2018towards}. The key challenge is the mapping from
symbolic prosody (the / and \textbackslash\ labels of SpeechPrint's tier) to
quantitative F0 targets that a synthesis system can use. One approach is to
learn this mapping from a corpus of annotated speech; another is to use a
parametric F0 model in which the symbolic labels directly control the slope,
height, and velocity of F0 movements.

A practical implementation using currently available tools would use the
Coqui TTS framework or a similar system that exposes duration and pitch control
parameters. Coqui XTTS, for example, allows voice cloning from a short reference
segment and accepts prosodic modifications through its Python API. The SpeechPrint
symbolic tier would be compiled into a sequence of pitch targets at word or
syllable boundaries, which the synthesis system would interpolate into a smooth
F0 contour. This would allow a researcher to hear the prosodic consequences of an
annotation change without recording a new utterance -- a capability that could be
valuable both for validation of the annotation and for generating stimuli in
prosody experiments.

\begin{figure}[htbp]
\centering
\fbox{\parbox{0.9\linewidth}{\centering\vspace{2mm}
\textbf{Analysis direction (current):} \\[2mm]
WAV $\rightarrow$ WhisperX $\rightarrow$ MFA / hardcoded IPA
$\rightarrow$ \{pYIN (read speech) $|$ CREPE (archival)\}
$\rightarrow$ Symbolic layer (/ \textbackslash\, * {\={}} \_)
$\rightarrow$ TextGrid \\[4mm]
\textbf{Resynthesis direction (proposed):} \\[2mm]
TextGrid (edited symbols) $\rightarrow$ F0 target compiler $\rightarrow$
Coqui TTS / XTTS $\rightarrow$ Resynthesized WAV \\[2mm]
}}
\caption{The proposed bidirectional SpeechPrint pipeline. The analysis direction
is fully implemented; the resynthesis direction is a design concept for future
work. The loop would allow a user to modify prosody symbols and hear the result
without re-recording.}
\label{fig:resynthesis_pipeline}
\end{figure}

\subsection{A Speech DAW: Interactive Prosody Manipulation}
\label{sec:speech_daw}

Taking the resynthesis concept further, a \emph{speech digital audio workstation}
(Speech DAW) would present the TextGrid annotation tiers as an editable timeline,
similar in concept to how a music DAW presents MIDI notes. The user would be able
to click on any syllable label in the prosody tier and change it, drag the onset
or offset of a pitch movement, raise or lower the height symbol from \texttt{\_}
to \texttt{{\={}}}, or add or remove an accent marker \texttt{*}. Each edit would
trigger a re-synthesis of the affected region, allowing real-time auditory
feedback.

The design is motivated by a gap in current tools. Praat allows the manual
manipulation of a pitch track at the acoustic level, which requires expert
knowledge of what a changed F0 contour sounds like. The proposed Speech DAW
would operate at the symbolic level, where the units (/ \textbackslash\, *) are
interpretable by a non-specialist. This lowers the barrier for users in
conversation analysis or field linguistics who want to test hypotheses about
prosodic structure by listening to counterfactual examples.

\subsection{Automatic Prosody Classification}
\label{sec:prosody_classifier}

A third future direction is the development of a classifier that assigns one of a
small number of prosodic categories to each sentence or utterance. The motivation
is that a single global label -- for example, whether a sentence is a statement, a
question, an emphatic assertion, a continuation, or a boundary -- is often more
useful for annotation workflows than a fine-grained per-syllable sequence of
symbols. The classification would be trained on existing annotated corpora and
applied to new recordings.

A candidate set of five prosodic categories, sufficient to cover the most
frequent intonational events in English and German, might be:
\begin{enumerate}
\item \textbf{Declarative}: broad-focus statement with final fall.
\item \textbf{Interrogative}: yes/no question with final rise.
\item \textbf{Narrow focus}: prominent accent on a single content word,
      marking contrast or correction.
\item \textbf{Continuation}: non-final intonational phrase with
      non-falling boundary tone, signalling that more speech follows.
\item \textbf{Exclamation}: emphatic or expressive contour, typically with
      high onset and rapid fall.
\end{enumerate}

This classification scheme is a simplification of the GToBI inventory but
captures the categories most relevant to the use cases described in
Chapter~\ref{scenarios}. The classifier would take the sequence of SpeechPrint
symbolic labels as input features, supplemented by the mean F0, F0 range, and
final boundary tone of the utterance. The output category would be displayed as
an additional tier in the TextGrid, providing a sentence-level interpretive
summary alongside the syllable-level labels.

\subsection{Bias, Projection, and Cross-Lingual Phonological Approximation}
\label{sec:bias_projection}

A recurring observation in this thesis is that applying a speech model trained
on one language to a recording in an unrelated language does not produce random
noise --- it produces something coherent but wrong. When Whisper's Italian
acoustic model was applied to the Daakie recording, it returned phonetically
plausible Italian word sequences, not silence or unintelligible garbage
(Appendix~\ref{app:daakie}). The model could not understand Daakie, but it
projected its learned phonological structure onto the signal and produced the
nearest thing in its training distribution. This is model bias operating as a
generative mechanism: the failure is systematic and, viewed from a different
angle, compositionally interesting. Each reinterpretation layer introduces a
transformation that is determined by what the model knows and what it does not,
producing emergent changes in apparent timbre, prosody, and intelligibility that
are not random but structured by the model's prior. The SpeechPrint prosody
layer applied to such a mistranscribed signal assigns labels that are acoustically
correct but linguistically misaligned, creating a kind of prosodic ghost of the
original utterance. This phenomenon --- the gap between what a model hears and
what was said --- represents a potential compositional space for work at the
intersection of speech technology and sound art, and connects to a broader
research thread on how bias and error in computational listening shape the
transformation of voice across systems.

% ENDED ADDING

\newpage
